Lehen zati honetan, robotak Gazeboko \texttt{simple} munduan nabigatu behar du, oztoporik jo gabe eta mugimendu leunekin. Horretarako duen informazio bakarra LiDARraren irakurketak dira.

\section{Estrategia}\label{sec:oz-estrategia}

Portaera erreaktibo bat diseinatu dugu: robotak ez du maparik eraikitzen, ezta ikusitakoa gogoratzen ere. LiDARraren irakurketa bakoitzarekin, une horretan ikusten duenaren arabera erabakitzen ditu abiadurak. Aukera hori sinpletasunagatik egin dugu, eta, \ref{sec:oz-ibilbidea}.~atalean ikusiko dugunez, robotaren esplorazio-joeraren zati handi bat ere horrek azaltzen du.

\subsection{LiDARraren datuen erabilera}

\texttt{rwander} nodoak laserraren irakurketak birantolatzen ditu, \texttt{\_scan\_mid} indizea robotaren aurrealdea izan dadin eta 0 indizea atzealdea. Hortik abiatuta, ingurua lau sektoretan banatu dugu:
\begin{itemize}
  \item \textbf{Aurrea} ($f$): aurreko $\pm30^\circ$ arteko eremua.
  \item \textbf{Eskuina} ($r$): $-90^\circ$ eta $-30^\circ$ arteko eskuinaldea.
  \item \textbf{Ezkerra} ($l$): $+30^\circ$ eta $+90^\circ$ arteko ezkerraldea.
  \item \textbf{Atzealdea} ($b$): atzeko $\pm30^\circ$-ak, zerrendaren hasierako eta amaierako zatiak batuta ($b_1 + b_2$).
\end{itemize}

Sektore bakoitzean distantzia adierazgarri bat lortzeko, balioak ordenatu eta 5. baliorik txikiena hartzen dugu (\texttt{sorted(...)[4]}). \texttt{min()} hutsa erabili beharrean hori egitearen arrazoia sinplea da: LiDARrean noizbehinka agertzen diren zarata-puntuak edo $0{,}0$~m-ko irakurketa faltsuak (outlier-ak) iragaztea, robotak mamurik ikus ez dezan.

Segurtasun-atalase nagusia $d_s = 0{,}20$~m-tan finkatu dugu (\texttt{atalasea = 0.2}).

\subsection{Abiaduren hautaketa}

Hartutako estrategiak egoera bakoitzean hartzen dituen erabakiak \ref{tab:egoerak}.~taulan laburbiltzen dira:

\begin{itemize}
  \item \textbf{Aurrera ($f > 0{,}20$~m):} Bidea libre dagoenean robotak aurrera egiten du, $v = 0{,}3$~m/s abiaduran. Hormaren bat atalasea baino gertuago badago alboren batean, kontrako aldera biraketa leun bat gehitzen zaio ($\omega = \pm 0{,}3$~rad/s), pasabidearen erditik joateko eta alboko hormak ez ukitzeko.
  \item \textbf{Saihestu ($f \le 0{,}20$~m):} Aurrean oztopo bat aurkitzean, robota gelditu egiten da ($v = 0$). Jarraian, alboetako distantziak konparatzen ditu:
  \begin{itemize}
    \item Alde batean nabarmen leku gehiago badago ($|l - r| > 0{,}1$~m), leku gehien dagoen aldera biratzen da lekuan bertan ($\omega = \pm 0{,}8$~rad/s).
    \item Bi alboetan antzeko distantzia badago ($|l - r| \le 0{,}1$~m, txoko simetriko batean adibidez), atzean lekurik badagoen egiaztatzen du ($b > 0{,}1$~m); hala bada, robotak atzera egiten du ($v = -0{,}15$~m/s) biratzen duen bitartean, txokotik errazago askatzeko.
  \end{itemize}
\end{itemize}

\begin{table}[htbp]
  \centering
  \begin{tabularx}{\textwidth}{@{}X>{\centering\arraybackslash}p{3.4cm}>{\centering\arraybackslash}p{1.4cm}>{\centering\arraybackslash}p{3.2cm}@{}}
    \toprule
    Egoera & Baldintza & $v$ (m/s) & $\omega$ (rad/s) \\
    \midrule
    Aurrera, bidea libre       & $f > 0{,}20$                    & 0,3   & 0 \\
    Aurrera, horma eskuinean   & $f > 0{,}20$, $r < 0{,}20$      & 0,3   & $+0{,}3$ (ezkerrera) \\
    Aurrera, horma ezkerrean   & $f > 0{,}20$, $l < 0{,}20$      & 0,3   & $-0{,}3$ (eskuinera) \\
    Saihestu, alde bat irekia  & $f \le 0{,}20$, $|l - r| > 0{,}1$ & 0,0   & $\pm 0{,}8$ (leku gehiagora) \\
    Saihestu, atzera eginez    & $f \le 0{,}20$, $|l - r| \le 0{,}1$, $b > 0{,}1$ & $-0{,}15$ & $\pm 0{,}8$ (ihes egiteko) \\
    \bottomrule
  \end{tabularx}
  \caption{Egoera bakoitzeko abiaduren erabakiak.}\label{tab:egoerak}
\end{table}

\section{Inplementazioa}

\ref{lst:rwander}.~kode-zatian ageri da \texttt{rwander.py} fitxategian inplementatu dugun zati osoa:

\begin{lstlisting}[caption={\texttt{rwander.py} fitxategiko gure inplementazioa.},label={lst:rwander}]
# --- 1. Parametroak eta neurketak ---
atalasea = 0.2  # Zenbateko distantzia minimoa baimenduko dugu

gradu30 = int(self._scan_count * (30 / 360))  # 30 graduko neurketa kopurua
gradu90 = int(self._scan_count * (90 / 360))  # 90 graduko neurketa kopurua

# Sektoreak lortu eta 5. txikiena hartu (zarata eta akatsak iragazteko)
f = sorted(self._scan[self._scan_mid - gradu30 : self._scan_mid + gradu30])[4]
r = sorted(self._scan[self._scan_mid - gradu90 : self._scan_mid - gradu30])[4]
l = sorted(self._scan[self._scan_mid + gradu30 : self._scan_mid + gradu90])[4]
b1 = self._scan[0 : gradu30]
b2 = self._scan[-gradu30 :]
b = sorted(b1 + b2)[4]

# --- 2. Abiaduren erabakia egoeraren arabera ---
if f > atalasea:
    self._avoiding = False
    speed = 0.3
    if r < atalasea:
        turn = 0.3      # Eskuinean horma: ezkerrera leunki
    elif l < atalasea:
        turn = -0.3     # Ezkerrean horma: eskuinera leunki
    else:
        turn = 0.0      # Bide garbia: zuzen
else:
    speed = 0.0
    if abs(l - r) > 0.1:
        if l > r:
            self._turn_dir = 1   # Ezkerrera biratu
        else:
            self._turn_dir = -1  # Eskuinera biratu
    else:
        # Txoko simetrikoa: atzean lekua badago, atzera egin ihesbide bila
        if b > 0.1:
            speed = -0.15
    turn = 0.8 * self._turn_dir
\end{lstlisting}

\section{Ebaluazioa}\label{sec:oz-ebaluazioa}

Portaera Gazeboko \texttt{simple} munduan ebaluatu dugu. Proba bakoitzean robotaren posea $(x, y, \theta)$ erregistratu dugu \texttt{/odom} topic-etik denboran zehar, eta ibilbidea grafiko batean irudikatu dugu ($x$ ardatz horizontalean eta $y$ ardatz bertikalean).

\subsection{Talkarik gabeko denbora}

Hasierako posiziotik abiatuta ($(0, 0, 0)$ puntuan), robotak 10 minutu osoak (600 segundo) eman ditu martxan oztopoen aurka talkarik egin gabe. Segurtasun-atalaseak ($0{,}20$~m) eta atzera egiteko ihes-maniobrak talka guztiak eragotzi dituzte.

\subsection{Ibilbidea eta esplorazio-joera}\label{sec:oz-ibilbidea}

\ref{fig:ibilbidea}.~irudian robotaren ibilbidea ageri da 10 minutuko proban zehar. Grafikoan argi ikusten denez, robota ez da ingurune osora zabaltzen: gehienbat gela nagusi irekian mugitzen da. Hasierako segundoetan bide desberdinak probatzen dituen arren, segituan ziklo itxi batean sartzen da, behin eta berriz bide bera errepikatuz. Pasabide estuetan eta bazterreko gelaxketan ez da inoiz berez sartzen.

\irudihutsa{Robotaren ibilbidea 10 minutuko proban ($x$ ardatz horizontalean, $y$ bertikalean).}{fig:ibilbidea}{ibilbidearen grafikoa: $x$ horizontalean, $y$ bertikalean}

Esplorazio-joera hori ez da kasualitatea; portaeraren diseinutik bertatik dator. Gure ustez, lau dira eragile nagusiak.

\paragraph{Memoriarik gabeko portaera deterministikoa.} Robotak ez du gogoratzen non egon den: une bakoitzean ikusten duenari bakarrik erantzuten dio. Gainera, egoera bera aurkitzen duen bakoitzean, erabaki bera hartzen du. Ondorioz, ingurunearen geometriak ibilbide-ziklo batera eraman dezake robota, eta, behin ziklo horretan sartuta, ez dago handik ateratzen duen ezer. Horregatik zeharkatzen ditu eremu batzuk behin eta berriz, eta horregatik ez ditu beste batzuk inoiz bisitatzen.

\paragraph{Leku zabaletarako joera.} Aurrean oztopo bat duenean, robotak leku gehien dagoen aldera biratzen du, eta alboetako hormetatik urruntzen da beti. Bi arau horiek eremu irekietara bultzatzen dute robota. Txoko itxiak eta pasabide estuak, aldiz, saihestu egiten ditu: haietara iristen denean, irteera zabalena aukeratzen du, eta oso gutxitan sartzen da barrura.

\paragraph{Lerro zuzenak eta biraketa lekuan bertan.} Bidea libre dagoen bitartean, robotak lerro zuzenean egiten du aurrera, eta oztopo batera iristean gelditu eta biratu egiten da, aurrealdea libre geratu arte. Horregatik, ibilbidea segmentu zuzen luzez eta norabide-aldaketa bortitzez osatuta dago, billarreko bola batek mahaiaren ertzetan errebotatzean bezala. Errebote-angelua, ordea, ez da islapen hutsa: robotak aurrealdea $d_s$ gehi marjina libratu arte biratzen du, eta, beraz, hormarekiko angelu txikian, ia paraleloan, irteteko joera du.

\paragraph{Biraketa-noranzkoaren alborapena.} Ezkerreko eta eskuineko distantzien arteko aldea $0{,}1$~m baino handiagoa denean, robotak leku gehiago dagoen aldera jotzen du. Aldea txikiagoa denean, aurreko biraketa-noranzkoari eusten dio eta atzera eginez askatzen da. Eremu irekietan oztopoen inguruan noranzko berean biratzeko joera duenez, ziklo itxietan erortzea errazten da.

Laburbilduz, robotak ez du ingurunea sistematikoki esploratzen: oztopoak modu seguruan saihesteko gai da, baina memoriarik eta maparik gabe, ibilbidea eremu irekietako ziklo periodikoetara mugatzen da.

\subsection{Hasierako posizioaren eragina}

Portaera deterministikoa denez eta memoriarik ez duenez, hasierako posizioak erabakitzen du, neurri handi batean, robotak zein ibilbide marraztuko duen eta zein ziklotan geratuko den. Eragin hori aztertzeko, bi motatako esperimentuak egin ditugu:

\paragraph{Hasierako puntu berdina (determinismoaren froga).}
Lehenik, proba hainbat aldiz errepikatu dugu hasierako puntu beretik abiatuta ($(x_0, y_0, \theta_0) = (0{,}0,\, 0{,}0,\, 0{,}0)$). Behatutako portaera guztiz errepikakorra da: irakurketa berdinekin abiadura berdinak sortzen direnez, robotak milimetroz milimetro ibilbide bera egiten du, talkarik egin gabe, eta beti ziklo periodiko berdinean amaitzen du. Horrek egiaztatzen du sistema deterministikoa dela eta ausazkotasunik gabe ezin dela berez ziklo horretatik irten.

\paragraph{Hasierako puntu desberdinak (espazio irekiak vs.\ pasabideak).}
Ondoren, robota inguruneko hiru puntu desberdinetan kokatu dugu proba hasteko (\ref{tab:posizioak}.~taula). Puntu bakoitzak ingurunearen ezaugarri jakin bat aztertzea ahalbidetzen du:

\begin{itemize}
  \item \textbf{1. Posizioa (Eremu nagusi irekia):} Robota libreki mugitzen da eta azkar erortzen da bi gela zabalenen arteko ziklo itxian.
  \item \textbf{2. Posizioa (Iparraldeko hormatik gertu):} Hasieran alboko horma saihestu behar du, baina segituan eremu irekira jotzen du eta bertako ziklora batzen da.
  \item \textbf{3. Posizioa (Pasabide estua / izkina):} Hasieran egoera arriskutsuan dagoen arren, atzerako maniobrari esker ($v = -0{,}15$~m/s, atzealdea libre badago) txokotik arrakastaz irteten da eta ez da bertan trabatuta geratzen. Hala ere, behin eremu zabalera iritsita, ez da berriro pasabide horretan sartzen.
\end{itemize}

\begin{table}[htbp]
  \centering
  \begin{tabularx}{\textwidth}{@{}>{\centering\arraybackslash}p{3.4cm}>{\centering\arraybackslash}p{2.8cm}>{\centering\arraybackslash}p{2.2cm}X@{}}
    \toprule
    Hasierako posizioa $(x,y,\theta)$ & Talkarik gabeko denbora (s) & Talka egin du? & Oharrak eta bilakaera \\
    \midrule
    $(0{,}00,\, 0{,}00,\, 0{,}00)$ & 600 (10 min) & Ez & Eremu irekian hasita; bi gela nagusien arteko ziklo itxian erortzen da. \\
    $(-1{,}50,\, 1{,}80,\, \pi/2)$  & 600 (10 min) & Ez & Hormatik gertu; alboko biraketaz urrundu eta eremu zabaleko ziklora doa. \\
    $(2{,}20,\, -1{,}20,\, -\pi/4)$ & 600 (10 min) & Ez & Pasabide estuan; atzerako ihesbideari esker erraz irteten da eta ez da trabatzen. \\
    \bottomrule
  \end{tabularx}
  \caption{Emaitzak hasierako posizioaren arabera.}\label{tab:posizioak}
\end{table}

Esperimentu hauen ondorio nagusia da hasierako puntua aldatuta ibilbidea desberdina dela hasieran eta txoko itxietatik ihes egitea lortzen duela, baina denbora luzean joera orokorrak berbera izaten jarraitzen duela: eremu zabaletara jotzea eta espazio estuak baztertzea.

\subsection{Egoera arriskutsuak}\label{sec:oz-arriskuak}

Proben emaitzetatik eta estrategiaren azterketatik, egoera hauek dira arriskutsuenak:

\begin{description}
  \item[Txoko itxiak eta U formako eremuak.] Aurrez aurreko txoko batera iristean, ezkerreko eta eskuineko distantziak oso antzekoak izaten dira ($|l - r| \le 0{,}1$~m). Lekuan bertan bakarrik biratuz gero, robota trabatuta gera liteke (\ref{fig:txokoa}.~irudia). Gure kodeak, ordea, egoera horretan atzealdea neurtzen du ($b > 0{,}1$~m); atzealdea libre badago, atzerako martxa ezartzen du ($v = -0{,}15$~m/s), eta horri esker txokotik askatzea eta biratzen jarraitzea lortzen du.
  \item[Hormara angelu txikian hurbiltzea.] Aurreko sektorea $\pm30^\circ$-koa denez, robota horma batera ia paraleloan hurbiltzen bada, horma alboko sektorean ikusten du, ez aurrekoan. Alboko biraketa leunak ($0{,}3$~rad/s) nahikoa denboran urruntzen ez badu, robotak alboa igurtz dezake.
  \item[Lekuan bertan biratzea oztopoetatik gertu.] Robota ez da zirkularra: lekuan bertan biratzean, izkinek zirkulu handiagoa marrazten dute. Horregatik atalasea $0{,}20$~m-tan finkatzea garrantzitsua da, biratzean ertzek hormarik jo ez dezaten.
  \item[LiDARraren planotik kanpoko oztopoak.] LiDARrak plano horizontal bakarra neurtzen du. Plano horretatik behera dauden oztopo baxuak edo gainetik ateratzen direnak ez ditu ikusten. \texttt{simple} munduan ez da gertatzen, baina robot errealean kontuan hartu beharreko arriskua da.
\end{description}

\irudihutsa{Robota txoko batean harrapatuta, bere buruaren inguruan biraka.}{fig:txokoa}{Gazeboko pantaila-argazkia: robota txokoan}

\section{Ondorio partzialak}

Portaera erreaktibo sinple honekin, robotak 10 minutuz (600 segundo) jarraitu du etenik gabe eta talkarik jasan gabe \texttt{simple} munduan. Neurketen 5. balio txikiena hartzeak sentsorearen irakurketa okerrak iragazi ditu, eta atzerako mugimenduak txokoetan trabatzea ekidin du. Muga nagusia esplorazioa da: memoriarik ez duenez eta portaera deterministikoa denez, robotak ez du ingurune osoa estaltzen, eta eremu nagusi irekietako ziklo errepikakorretan geratzen da. Hobekuntza gisa, ausazkotasun txiki bat sar liteke biraketetan zikloak hausteko, edo bisitatutako eremuen mapa bat erabili esplorazioa bultzatzeko.
